\section{Introduction}
\label{sec:introduction}

\begin{sloppypar}
The structural assessment of Unreinforced Masonry (URM) buildings following catastrophic events, such as seismic activity or urban warfare, is a fundamental pillar of disaster resilience and humanitarian recovery \cite{cervenka2020computer}. URM structures, characterized by their high mass and low tensile capacity, remain one of the most vulnerable construction typologies globally. Recent seismic events in high-seismicity regions (e.g., the 2015 Gorkha earthquake in Nepal, the 2023 Turkey-Syria earthquake sequence) have underscored the critical need for rapid, scalable, and engineering-grade damage evaluation. Traditionally, this process relies on "Rapid Visual Screening" (RVS) performed by human inspectors. However, manual inspections are inherently subjective, time-intensive, and often lack the quantitative grounding required for deterministic "safe/unsafe" occupancy labeling. 
\end{sloppypar}

\subsection{Evolution of Computer Vision in Structural Inspection}
\begin{sloppypar}
The advent of Deep Learning (DL) has catalyzed a paradigm shift in structural health monitoring (SHM). Convolutional Neural Networks (CNNs) and transformer-based architectures, such as YOLOv8 and ResNet, have demonstrated exceptional performance in localized crack detection and multi-class damage classification. Despite these advances, a significant "diagnostic gap" persists. While state-of-the-art models can identify pixel-level fissures with high precision, they typically lack the semantic context required to interpret the \textit{structural significance} of a crack. In the domain of professional structural engineering—specifically as defined by standards such as \textbf{FEMA 306}—a crack is not merely a geometric abnormality; it is a manifestation of a specific failure mechanism (e.g., Toe Crushing, Bed Joint Sliding, or Diagonal Tension) determined by the crack's morphology, its location within the wall assembly (pier vs. spandrel), and its propagation history.
\end{sloppypar}

\subsection{The Semantic Challenge: Pixel-Level vs. Engineering-Grade Diagnosis}
\begin{sloppypar}
Existing automated systems often treat crack detection as a generic object detection task, ignoring the procedural decomposition of the wall itself. A $5mm$ diagonal crack in the center of a pier indicates a fundamentally different structural risk than a similar crack at a spandrel-pier intersection. Bridging this gap requires a multi-modal synthesis of three distinct data streams: (1) high-fidelity morphological features (width, length, orientation), (2) architectural context (wall component partitioning), and (3) codified engineering knowledge (RAG-based inference). 
\end{sloppypar}

\subsection{Research Contributions and Framework Overview}
\begin{sloppypar}
This paper proposes an integrated, hierarchical framework for the autonomous assessment of masonry walls, transitioning from raw imagery to auditable engineering reports. The core contributions of this research are fourfold:
\end{sloppypar}

\begin{enumerate}[noitemsep]
    \item \textbf{Hierarchical Multi-Stage CV Pipeline:} We implement a tiered approach using YOLOv8 for rapid localization and a GAN-refined DeepCrack architecture for pixel-level morphological extraction, ensuring robust performance across varying textures and scales.
    \item \textbf{Autonomous Scale Calibration via Pattern Recognition:} We introduce a novel calibration logic that leverages the periodic nature of brick courses (via FFT and morphological projection) to establish a pixel-to-mm scale factor without the need for external markers.
    \item \textbf{Vision-Language Wall Decomposition:} Utilizing zero-shot Vision-Language Models (VLMs) and procedural computer vision, we decompose wall structures into their constituent piers and spandrels, providing the "topological anchor" needed for structural mapping.
    \item \textbf{Knowledge-Based RAG Inferece with Hybrid Scoring:} We propose a two-layer RAG architecture that queries digitized FEMA 306 standards to generate diagnostic reports. To ensure reliability, we introduce a \textbf{Hybrid Confidence Scoring (HCS)} mechanism that ensembles LLM self-confidence, retrieval distance, and cross-encoder re-ranking.
\end{enumerate}

\begin{sloppypar}
The remainder of this paper is organized as follows: Section \ref{sec:literature_review} provides a comprehensive review of existing methodologies; Section \ref{sec:materials_methods} details the proposed hierarchical framework, including data collection and preprocessing; Section \ref{sec:results} presents the experimental validation and detailed analysis of the pipeline's performance; Section \ref{sec:discussion} discusses the implications and limitations of the study; and Section \ref{sec:conclusion} concludes the paper with future research directions.
\end{sloppypar}
